\chapter{Generátory grafů \texttt{geng} a \texttt{plantri}}
\label{chap:generators}

V předchozí kapitole jsme popsali obecný problém generování grafů bez izomorfních
kopií. V této kapitole se podíváme na dva konkrétní generátory, na kterých je tato práce
postavena: \texttt{geng} z balíku nauty a \texttt{plantri}. Oba nástroje řeší podobný
typ úlohy, ale pracují s jinými třídami grafů a používají jiné interní reprezentace.
Pro návrh jednotného rozhraní proto potřebujeme porozumět nejen tomu, jaké grafy
generují, ale také tomu, jak se do jejich běhu dá zasáhnout z vlastního programu.

\section{Generátor \texttt{geng}}
\label{sec:geng}

Generátor \texttt{geng} je součástí balíku nauty. Balík nauty je známý především
algoritmy pro práci s automorfismy grafů a s kanonickým značením~\cite{mckay2014practical}.
Program \texttt{geng} využívá tyto algoritmy pro generování neizomorfních neorientovaných
grafů na zadaném počtu vrcholů. Uživatel může omezit například počet hran, souvislost,
minimální a maximální stupeň vrcholů nebo zakázat vybrané malé podgrafy
(například trojúhelníky či čtyřcykly).

Takový graf můžeme popsat množinou vrcholů a množinou hran. Při izomorfismu nás
tedy zajímá pouze zachování sousednosti. Reprezentace grafu v nauty je proto
založena na bitových množinách
sousedů, které jsou velmi kompaktní a dobře se hodí pro rychlé množinové operace.

Z pohledu uživatele na příkazové řádce je práce s \texttt{geng} poměrně přímočará:
zadáme počet vrcholů a přepínače popisující požadovanou třídu grafů. Uvnitř
však \texttt{geng} generuje grafy inkrementálně: z grafu na
\(n\) vrcholech vytváří graf na \(n+1\) vrcholech tak, že přidá nový poslední
vrchol a zvolí jeho sousedy mezi již existujícími vrcholy. Aby přitom nevznikaly
izomorfní kopie téhož grafu, nepřijme každé možné rozšíření. Nově vzniklý graf
musí projít kanonickým testem, který mimo jiné rozhoduje, zda právě poslední
přidaný vrchol může být považován za kanonicky zvolený.
Při tomto rozhodování se používají i stupně vrcholů. Zjednodušeně řečeno,
generátor nepracuje se všemi pořadími vzniku vrcholů stejně, ale využívá
vlastnosti posledního vrcholu k tomu, aby omezil počet zkoušených možností.

Při rozšíření generátoru o barvy je důležité, že barvy musí ovlivňovat také
testy izomorfismu a kanonické rozhodování během generování. Konkrétní úpravy
potřebné pro tuto podporu popíšeme v kapitole o implementaci.


\section{Generátor \texttt{plantri}}
\label{sec:plantri}

Generátor \texttt{plantri} je specializovaný nástroj pro generování rovinných grafů.
Umí generovat například rovinné triangulace, kvadrangulace, jednoduché rovinné grafy,
triangulace disku. Stejně jako \texttt{geng} negeneruje všechny možné označené grafy, 
ale jeden graf pro každou izomorfní třídu.
%Principy fungování nástroje \texttt{plantri} popisují Brinkmann a McKay v článku
%o rychlém generování rovinných grafů~\cite{brinkmann2007fast}.

Generování v \texttt{plantri} je založeno na postupném zvětšování rovinného
grafu pomocí lokálních operací, které se nazývají expanze. Opačné operace,
které graf zmenšují, se nazývají redukce~\cite[Section~1.2]{brinkmann2007fast}.
Jednoduše řečeno, expanze je krok, kterým graf zvětšíme, zatímco redukce je
odpovídající krok, kterým jej zmenšíme.

Například při generování triangulací můžeme začít od malého počátečního grafu a
postupně jej zvětšovat operacemi, které přidávají vrchol nebo upravují lokální okolí
některé stěny. Každé takové rozšíření zachovává rovinné vložení a strukturu
triangulace. Protože stejný výsledný graf může vzniknout více různými posloupnostmi
expanzí, musí generátor průběžně odmítat nekanonické větve výpočtu.

Tento problém řeší metoda kanonické konstrukční cesty~\cite[Section~1.2]{brinkmann2007fast}.
Pro nově vzniklý graf se určí kanonicky zvolená redukce, tedy opačný krok, kterým by
se graf měl zmenšit. Rozšíření se přijme pouze tehdy, pokud právě provedená expanze
odpovídá této kanonické redukci. Díky tomu generátor vytvoří pouze jeden reprezentant
z každé izomorfní třídy.


\texttt{plantri} generuje rovinné grafy spolu s jejich vložením. U takového
grafu nestačí vědět, které dvojice vrcholů jsou spojeny hranou. Potřebujeme
také znát cyklické pořadí hran kolem každého vrcholu, protože právě toto pořadí
určuje stěny rovinného vložení. \texttt{plantri} proto používá reprezentaci založenou na orientovaných
hranách. Každá neorientovaná hrana je reprezentována dvojicí opačně orientovaných hran.
Každá taková orientovaná hrana obsahuje odkazy na následující a předchozí hranu
v cyklickém pořadí kolem vrcholu.

Taková half-edge reprezentace umožňuje přirozeně procházet okolí vrcholu i stěny
rovinného vložení.


\section{Možnosti rozšíření pomocí callbacků}
\label{sec:generator-callbacks}

Oba generátory umožňují zasáhnout do generování vlastním kódem, ale dělají to způsobem
typickým pro C programy. \texttt{geng} nabízí například mechanismy \texttt{OUTPROC},
\texttt{PRUNE} a \texttt{PREPRUNE}. \texttt{OUTPROC} umožňuje nahradit standardní výpis
grafu vlastní funkcí. \texttt{PRUNE} a \texttt{PREPRUNE} umožňují odmítat průběžně vytvářené grafy
a tím prořezávat prohledávací strom.

\texttt{plantri} používá podobnou myšlenku přes pluginy a makra jako \texttt{FILTER}
a \texttt{PRE\_FILTER}. Uživatel může doplnit vlastní testy, které se spouštějí při
generování, a rozhodovat, zda má být graf vypsán nebo zda má být větev výpočtu
zahozena. Díky tomu můžeme ušetřit více času než při filtrování hotového 
výstupu, protože nevhodné části prohledávacího prostoru zahodíme dříve.

Tyto mechanismy jsou výkonné, ale jejich přímé použití z C++ vyžaduje znalost
přesných signatur callbacků, způsobu sestavení generátoru s příslušnými makry
a často také interní reprezentace aktuálního grafu.

\section{Omezení původních rozhraní}
\label{sec:original-interface-limitations}

Původní rozhraní generátorů \texttt{geng} a \texttt{plantri} je velmi vhodné pro použití
z příkazové řádky. Pokud chceme pouze vygenerovat grafy dané třídy a uložit je do
souboru, jsou oba nástroje rychlé a praktické. Pro použití jako součást většího C++
programu však narážíme na několik omezení.

Prvním omezením je rozdílný způsob konfigurace. Každý generátor má vlastní sadu
přepínačů, vlastní výstupní formáty a vlastní pojmenování podobných voleb. Druhým
omezením je nízkoúrovňové C rozhraní. Pokud chceme generátor volat přímo z programu,
musíme připravovat argumenty ve tvaru podobném příkazové řádce a pracovat s globálním
stavem původního generátoru.

Třetím omezením je reprezentace grafů v callbackách. U \texttt{geng} pracujeme s
bitovou reprezentací nauty, u \texttt{plantri} s poli stupňů a ukazateli na orientované
hrany. Obě reprezentace jsou efektivní, ale uživatel musí znát jejich detaily, pokud chce
nad grafem spustit vlastní algoritmus. Převod každého grafu do jiné datové struktury by
rozhraní zjednodušil, ale při velkém počtu generovaných grafů by mohl znamenat
výraznou ztrátu výkonu.

Při návrhu rozhraní proto chceme tato omezení skrýt, aniž bychom ztratili
výhody původních generátorů.
